% =======================================================================
\section{Discussion}

\subsection{Operating range}

The result that matters for use is the width of the operating range.
$H_1$ detection holds across cutoffs from 6 to 14~m. It also holds
across every truncation percentile from $P_{50}$ to $P_{95}$ tested
(Section~\ref{sec:sensitivity}). A pipeline that works only at one
finely tuned setting would not travel. This one does, once interaction
lengths are re-derived on a new domain.

The three $H_0$ regimes (individual, tactical, and team) are the
cardinality targets of Section~\ref{sec:cutoffselection}. The metres
that select them are read from the pooled
\texorpdfstring{$H_0(\delta)$}{} curve on every complete frame of the
ten matches (Figure~\ref{fig:cutoff}): $2.75$, $11.75$, and $23.0$~m.
Quality metrics are a check. They are not the selector.
Section~\ref{sec:sensitivity} shows that tactical $H_1$ is not a
spike at $11.75$~m.

$P_{75}$ in equation~\eqref{eq:adaptive} is a reporting convention.
All tested percentiles return identical $H_1$ totals and presence
rates.

\subsection{Distinct structure at two levels}

The weak correlation and the divergent presence rates in
Section~\ref{sec:complementarity} support the same conclusion. The two
levels are not measuring the same structure at different resolutions.
If the tactical signal were a coarsened version of the individual
signal, the two would co-occur far more consistently than
Table~\ref{tab:fisher} shows. The bottleneck and landscape distances
would also be small relative to each level's own persistence values.
Clustering before persistent homology is therefore not a convenience
for separating noise. It produces two informationally distinct
objects.

\subsection{Limitations}
\label{sec:limitations}

Linkage selection is a substantive methodological choice. Across a
600-frame, four-match comparison sample, single-linkage detects 106
tactical-level $H_1$ loops. Complete-linkage detects 802. Ward's
method detects 822. That is a difference of almost an order of
magnitude. The gap is large enough to change which scientific picture
the data support, not merely the precision of one estimate. We retain
single-linkage because its definition is the most natural reading of
spatial proximity. It is also robust against over-partitioning
artefacts. We treat the gap with the other two methods as an open
question.

A more principled selection criterion would come from the dynamical
system itself, rather than from clustering-quality metrics alone. One
candidate is the linkage under which a governing equation for a
persistence functional admits the sparsest representation, in the
sense of sparse identification of nonlinear dynamics. The functional
would be the landscape $L^2$ norm or the tactical $H_1$
total-persistence sequence. Operationalising this would require three
things. First, choose the dynamical state. Second, defend the sparsity
of the recovered library against AIC, BIC, or cross-validation
baselines. Third, confirm that the resulting persistence sequence is
genuinely lower-dimensional under the selected linkage. We defer this
to the forthcoming full-season work. A population-sized sample of
match sequences is what makes the sparsity comparison meaningful.

The tactical cutoff is the inversion of a named cardinality (mean
$H_0$ nearest 5, with $k\ge 4$ still common). On the 1~Hz diagnostic
subset the only interior silhouette local maximum is at $7.0$~m, a
small-group alternative. Silhouette is omitted at $k=1$. Among frames
that still have $k\ge 2$ it then keeps rising into the two-envelope
regime. Those features are alternatives, not adopted cutoffs. A
landscape-path stability criterion (minimise variation of the
persistence landscape under small perturbations of $\delta$) remains a
future estimator. It is computable with the landscape infrastructure of
Section~\ref{sec:complementarity}. We defer it to the
persistence-landscape companion work.
Section~\ref{sec:sensitivity} reports the operative tactical $H_1$
range as an empirical property around the adopted $11.75$~m.

The pipeline identifies three $H_0$ regimes but only two $H_1$
regimes in the headline analysis. At $\delta=23.0$~m the cloud has
mean $H_0=1.98$ and $k\le 2$ in 93.6\% of complete frames. A
Vietoris--Rips complex on at most two points cannot carry a
non-trivial 1-cycle. Residual frames with $k=3$ are few (6.4\%).
Team-level loop structure on the typical envelope would require a
different representation, such as spatial density fields or Delaunay
triangulations.

All data analysed here are broadcast-derived tracking at 10~Hz.
Broadcast tracking typically offers lower spatial precision than
higher-frequency optical systems. That may affect the magnitude of
persistence values reported. The structural findings (the three $H_0$
regimes, the $H_1$ presence rates, and the sensitivity profiles)
depend on relative rather than absolute spatial precision. They are
therefore expected to be robust across tracking technologies. A
direct comparison of persistence magnitudes against optical tracking
data remains for future work.

\subsection{Outlook}

Two extensions are in progress. A third output sits alongside them
rather than ahead of them. The ten-match evidence base is being scaled
to a full season of Championship matches. That sample will
characterise population-level distributions of $H_0$ and $H_1$
counts, barcode lengths, and landscape norms, and will test
tactical-fingerprint classification at population scale. It is also
the setting in which the linkage criterion proposed above becomes
evaluable. Recovering a sparse governing equation for a persistence
functional requires the population-sized sample of match sequences
that a single season provides. Separately, persistence landscape
dynamics are being developed to treat each match as a path through
landscape space. That is the setting in which the landscape-stability
cutoff criterion can itself be evaluated. Its resampling unit is the
match-level landscape, not the per-frame summary used here.

A companion paper \citep{paperB} interprets the present ten-match
findings for a football-analytics readership. Event correlation is a
central result there rather than a brief validity check. That paper
also compares the pipeline against standard geometric descriptors,
examines bilateral home-and-away coupling, and tests predictive
utility for phase-of-play classification. It depends methodologically
on the pipeline validated here. It asks what these topological
measures reveal about football specifically, rather than whether the
pipeline itself is sound.
